%\subsubsection{Verifica dei documenti}
%Ogni documento prodotto dal team passa per una o più fasi di verifica prima di essere pubblicato ufficialmente nel repository \textit{GitHub}.\\
%Tale fasi vengono realizzate solamente in seguito alla consegna del documento da parte dei responsabili di stesura, i quali ritengono che in quel momento il prodotto da loro realizzato sia completato.\\
%La verifica dei documenti prevede che venga assicurato:
%\begin{itemize}
%    \item Il rispetto delle norme tipografiche stabilite nelle \textit{Norme di progetto} del team;
%    \item Il rispetto della struttura definita nelle \textit{Norme di progetto} del team;
%    \item Il rispetto del contenuto, il quale deve essere chiaro, coerente, preciso ed esaustivo;
%    \item Il rispetto della sintassi italiana e della correttezza ortografica dei termini utilizzati.
%\end{itemize}

\documentclass[10pt, a4paper]{article}

% Parametri che modificano il file main.tex
% Le uniche parti da cambiare su main.tex sono:
% - tabella versioni
% - testo file

\def\titolo{Analisi dei \\ \ requisiti} % \\ per andare a capo, \ per spazio
\def\titoloHeader{Analisi dei requisiti}
\def\data{v2.0.0(0)}

% \def\listaComponenti{
% Bresolin G.,
% Campese M.,
% Ciriolo I.,
% Dugo A.,
% Feltrin E.,
% Michelon R.,
% Orlandi G.
% }

\usepackage{style}
\usepackage{headerfooter}
\usepackage{comment}

\title{\titolo}
\author{SWEetCode}

\begin{document}

% PRIMA PAGINA
\begin{titlepage}
    \thispagestyle{empty}
    \begin{tikzpicture}[remember picture, overlay]
        % TRIANGOLI
        \draw[fill=secondarycolor, secondarycolor] (current page.north west) -- (current page.south west) -- (8.8, -28);
        \draw[fill=primarycolor, primarycolor] (-3, 5) -- (4, -13.6) -- (11, 5);

        % LOGO
        \node [xshift=-5cm, yshift=25cm] (logo) at (current page.south east) {\includegraphics[width=6.5cm]{img/logo.png}};

        % SWEETCODE - DATE
        \node [anchor=north east, align=right, xshift=-1.2cm, yshift=20.5cm, text=black] (sweetcode) at (current page.south east) {\fontsize{32pt}{36pt}\selectfont SWEetCode};
        \draw[line width=4pt, lightcol] ([xshift=-3cm, yshift=-0.37cm]sweetcode.south west) -- ([yshift=-0.37cm]sweetcode.south east);
        \node [anchor=north east, align=right, xshift=-1.2cm, yshift=18.7cm, text=black] (date) at (current page.south east){\fontsize{24pt}{24pt} \selectfont Verbale \tipoVerb};

        % NOME FILE
        \node [anchor=north east, text width=15cm, align=right, xshift=-1.2cm, yshift=17cm, text=black] (titolo) at (current page.south east){\fontsize{48pt}{48pt}\textbf{\data}};

        % BOX DATI PARTECIPANTI
        \ifthenelse{\equal{\tipoVerb}{Esterno}}{
            \node[anchor=north east, xshift=-1.2cm, yshift=14.5cm, minimum width=8cm] (box) at (current page.south east){};
        }{
            \node[anchor=north east, xshift=-1.2cm, yshift=12.5cm, minimum width=8cm] (box) at (current page.south east){};
        }

        % RESPONSABILE
        \node[anchor=north west, align=left] (dati2) at (box.north west) {\fontsize{15pt}{15pt}\selectfont \textbf{Responsabile}};
        \draw[line width=4pt, lightcol] (dati2.south west) -- ([xshift=8cm]dati2.south west);
        \node[anchor=north west, align=left] (dati21) at (dati2.south west){\fontsize{13pt}{13pt}\selectfont \nomeResp};

        % VERIFICATORE
        \node[anchor=north west, yshift=-1cm, align=left] (dati3) at (dati21.north west) {\fontsize{15pt}{15pt}\selectfont \textbf{Verificatore}};
        \draw[line width=4pt, lightcol] (dati3.south west) -- ([xshift=8cm]dati3.south west);
        \node[anchor=north west, align=left] (dati31) at (dati3.south west){\fontsize{13pt}{13pt}\selectfont \nomeVer};

        % SEGRETARIO DI RIUNIONE
        \node[anchor=north west, yshift=-1cm, align=left] (dati4) at (dati31.north west) {\fontsize{15pt}{15pt}\selectfont \textbf{Segretario di Riunione}};
        \draw[line width=4pt, lightcol] (dati4.south west) -- ([xshift=8cm]dati4.south west);
        \node[anchor=north west, align=left] (dati41) at (dati4.south west){\fontsize{13pt}{13pt}\selectfont \nomeSegr};
        
        % UNIPD - SWE
        \node [xshift=4.4cm, yshift=2.3cm, draw, secondarycolor, text=white] (uni) at (current page.south west) {\fontsize{20pt}{20pt} \selectfont Università di Padova};
        \node [xshift=0.65cm, yshift=0.7cm, draw, secondarycolor, text=white, below=of uni] (corso) {\fontsize{20pt}{20pt}\selectfont Ingegneria del Software};

        % FIRMA AZIENDA
        \ifthenelse{\equal{\tipoVerb}{Esterno}}{
            \draw[line width=4pt, lightcol] ([xshift=-1.2cm, yshift=7cm]current page.south east) -- ([xshift=-8cm, yshift=7cm]current page.south east);
            \node [xshift=-4.8cm, yshift=8cm] (logo) at (current page.south east) {\includegraphics[width=6cm]{img/firme/\firmaAzienda}};
            \node[anchor=north west, xshift=12.9cm, yshift=6.65cm, align=left] at (current page.south west)
            {\fontsize{13pt}{13pt}\selectfont L'azienda: \nomeAzienda};
        }

        % FIRMA CODING COWBOYS
        \draw[line width=4pt, lightcol] ([xshift=-1.2cm, yshift=4.35cm]current page.south east) -- ([xshift=-8cm, yshift=4.35cm]current page.south east);
        \node [xshift=-4.8cm, yshift=5.5cm] (logo) at (current page.south east) {\includegraphics[width=6cm]{img/firme/menon.png}};
        \node[anchor=north west, xshift=12.9cm, yshift=4cm, align=left] at (current page.south west)
        {\fontsize{13pt}{13pt}\selectfont Coding Cowboys: Menon G.};

        % FIRMA
        \draw[line width=4pt, lightcol] ([xshift=-1.2cm, yshift=1.8cm]current page.south east) -- ([xshift=-8cm, yshift=1.8cm]current page.south east);
        \node [xshift=-4.8cm, yshift=2.45cm] (logo) at (current page.south east) {\includegraphics[width=6cm]{img/firme/\firmaResp}};
        \node[anchor=north west, xshift=12.9cm, yshift=1.45cm, align=left] at (current page.south west)
        {\fontsize{13pt}{13pt}\selectfont Il Responsabile: \nomeResp};
        
    \end{tikzpicture}
\end{titlepage}

% REGISTRO DELLE VERSIONI
%Registro in ordine dalla più recente alla meno recente!
{\renewcommand{\arraystretch}{1.5}
\section*{Registro delle versioni}

\begin{xltabular}{\textwidth}{c|c|c|c|X}
\label{tab:long}

\textbf{Versione} & \textbf{Data} & \quantities{\textbf{Responsabile di}\\\textbf{stesura}}& \textbf{Revisore} & \quantities{\textbf{Dettaglio e}\\\textbf{motivazioni}} \\
\endfirsthead

\textbf{Versione} & \textbf{Data} & \quantities{\textbf{Responsabile di}\\\textbf{stesura}}& \textbf{Revisore} & \quantities{\textbf{Dettaglio e}\\\textbf{motivazioni}} \\
\endhead

\multicolumn{5}{r}{{Continua nella pagina successiva}} \\
\endfoot

\endlastfoot

\hline
v1.10.0(7) & $2023-12-29$ & \quantities{Michelon R.} & Dugo A. & Aggiunti diagrammi use case e di attività.\\
\hline
v1.10.0(6) & $2023-12-29$ & \quantities{Michelon R.} & Dugo A. & Aggiornamento use case e requisiti.\\
\hline
v1.0.0(14) & $2023-12-06$ & \quantities{Bresolin G.} & Orlandi G. & Aggiornamento UC e requisiti presenti. Introduzione diagrammi dei casi d'uso. Aggiunta di nuovi UC e nuovi requisiti.\\
\hline
v1.0.0(5) & $2023-11-23$ & \quantities{Ciriolo I.} & Campese M. & Introduzione, Descrizione, Studio dei primi Casi d'uso, Studio dei primi requisiti.\\

    
\end{xltabular}}
\newpage

% INDICE
\tableofcontents
\newpage
\listoffigures
\newpage
\listoftables
\newpage

% INTRODUZIONE
\section{Introduzione}
\subsection{Scopo del documento}
L'obiettivo che ci si pone nella realizzazione di questo documento è la definizione delle metriche di valutazione e validazione del progetto, e la specifica degli obiettivi di qualità del prodotto finale. Tali parametri vengono stabiliti in accordo ai requisiti e 
alle aspettative del proponente e talvolta a discrezione del team sulla base delle valutazioni fatte nel corso di studi.

\subsection{Glossario}
Al fine di evitare incomprensioni relative alla terminologia usata all’interno del documento viene fornito un Glossario nel file apposito. Ogni termine presente in tale documento presenta una scrittura in corsivo ed un pedice |g|.

\subsection{Riferimenti}
   \subsubsection{Riferimenti normativi}
   \begin{itemize}
    \item \textit{Norme di progetto};
    \item \textit{Regolamento del progetto didattico}: \\
    \href{https://www.math.unipd.it/~tullio/IS-1/2023/Dispense/PD2.pdf}{https://www.math.unipd.it/~tullio/IS-1/2023/Dispense/PD2.pdf};
    \item \textit{Guidelines for
    the application of ISO 9001:2000 to
    computer software}: \\
    \href{https://cdn.standards.iteh.ai/samples/35867/36860aa4caba4c84b26051db576456d3/ISO-IEC-90003-2004.pdf}{https://cdn.standards.iteh.ai/samples/35867/36860aa4caba4c84b26051db576456d3/ISO-IEC-90003-2004.pdf}.
    \end{itemize}
    
    \subsubsection{Riferimenti informativi}
    \begin{itemize}
    \item \textit{Presentazione capitolato C1}: \\
    \href{https://www.math.unipd.it/~tullio/IS-1/2023/Progetto/C1p.pdf}{https://www.math.unipd.it/~tullio/IS-1/2023/Progetto/C1p.pdf};
    \item \textit{Analisi dei requisiti};
    \item \textit{Piano di progetto};
    \item \textit{Verbali esterni ed interni};
    \item \textit{Qualità di processo (argomento T8)}: \\
    \href{https://www.math.unipd.it/~tullio/IS-1/2023/Dispense/T8.pdf}{https://www.math.unipd.it/~tullio/IS-1/2023/Dispense/T8.pdf};
    \item \textit{Verifica e Validazione: Introduzione (argomento T9)}: \\
    \href{https://www.math.unipd.it/~tullio/IS-1/2023/Dispense/T9.pdf}{https://www.math.unipd.it/~tullio/IS-1/2023/Dispense/T9.pdf}.
    
   
    \end{itemize}

    
% Obiettivi di qualità
\newpage
\section{Obiettivi di qualità}
Per garantire la qualità di processo si è deciso di aderire allo standard ISO/IEC 90003:2004. In questa sezione si presentano i valori accettabili e ideali delle metriche che si rifanno a questo standard; le definizioni di tali metriche sono riportate nella sezione apposita (\color{red}6.3.4 Lista delle metriche\color{black}) del documento Norme di progetto(\color{red}Aggiungere versione\color{black}).
\subsection{Qualità di processo}
%\paragraph{}breve descrizione qualità di processo da inserire.\\
 {\renewcommand{\arraystretch}{1.5}
 \begin{table}[H]
    \begin{tabularx}{\textwidth}{p{0.1\textwidth}|p{0.5\textwidth}|X|X}
    \textbf{ID} & \textbf{Nome metrica} & \textbf{Valore tollerabile} & \textbf{Valore ambito}  \\
    \hline
     M1PMS & Percentuale di metriche soddisfatte & $ \ge75\% $& $\ge95\% $\\
    \hline
     M2RNP & Rischi non previsti & $ \le2 $& $\le1 $\\
    \hline
     M3VP & Variazione di piano & $ \ge-8 $ & $ \ge0 $ \\
    \hline
     M4VC & Variazione di costo & $0$ & $\le0 $ \\
    \hline
     M5ISR &  Indice di stabilità dei requisiti & $ \ge70\% $ & $ 100\% $  \\
    \hline
     M6CCM & Complessità ciclomatica media & $\le6 $ & $\le3 $ \\
    \hline
    M7SC & Statement coverage & $ \ge75\% $ & $ 100\% $ \\
    
    \hline
     M8BC & Branch Coverage & $ \ge75\% $ & $ 100\% $ \\
    \end{tabularx}
    \caption{Metriche per la qualità di processo}
    \end{table}
    
\subsection{Qualità di prodotto}
Abbiamo analizzato quali caratteristiche fossero necessarie per la realizzazione di un prodotto di 
qualità.\\
 {\renewcommand{\arraystretch}{1.5}
 \begin{table}[H]
  \begin{tabularx}{\textwidth}{p{0.1\textwidth}|p{0.5\textwidth}|X|X}
    \textbf{ID} & \textbf{Nome metrica} & \textbf{Valore tollerabile} & \textbf{Valore ambito}  \\
    %\endfirsthead
    \hline
     M9IG & Indice di Gulpease & $ \ge40\% $& $\ge80\% $\\
    \hline
     M10PRI & Percentuale di requisiti implementati & $ \ge85\% $ & $ 100\% $ \\
    \hline
     M11LMC & Linee medie di codice per metodo & $\le30$ & $\le20$ \\
    \hline
    M12DC & Densità dei commenti & $\ge25\%$ & $\ge30\%$ \\
    \hline
     M13TUS & Percentuale di test di unità superati & $ \ge85\% $& $100\%$\\
    \hline
     M14TIS & Percentuale di test di integrazione superati & $ \ge85\% $& $100\%$\\
    
    \end{tabularx}
    \caption{Metriche per la qualità di prodotto}
    \end{table}

    
%\subsection{Software}





\subsection{Qualità per obiettivo}
\subsubsection{Processi primari}

 
    {\renewcommand{\arraystretch}{1.5}
    \begin{tabularx}{\textwidth}{p{0.05\textwidth}|p{0.35\textwidth}|X|X}
    \textbf{ID} & \textbf{Nome metrica} & \textbf{Valore accettabile} & \textbf{Valore ottimale}  \\
    \hline
    \multicolumn{4}{l}{\cellcolor{primarycolor}\textbf{\textit{Fornitura}}} \\
    \hline
     &  &  &  \\
    \hline
     &  &  &  \\
    \hline
     &  &  &  \\
    \hline
     &  &  &  \\
    \hline
     &  &  &  \\
    \hline
     &  &  &  \\
    \hline
     &  &  &  \\
    \hline
    \multicolumn{4}{l}{\cellcolor{primarycolor}\textbf{\textit{Sviluppo}}} \\
    \hline
     &  &  &  \\
    \hline
     &  &  &  \\
    \end{tabularx}}

% PROCESSI DI SUPPORTO
\subsubsection{Processi di supporto}
\begin{comment}
%\subsubsection{Documentazione}
    {\renewcommand{\arraystretch}{1.5}
    \begin{tabularx}{\textwidth}{p{0.18\textwidth}|p{0.6\textwidth}|X}
    \textbf{Obiettivo} & \textbf{Descrizione} & \textbf{Metriche}  \\
    \hline
    Correttezza Risposte &   & \\
    \hline
    Velocità di esecuzione &  &  \\
    \end{tabularx}}
    {\renewcommand{\arraystretch}{1.5}
    \begin{tabularx}{\textwidth}{p{0.18\textwidth}|p{0.6\textwidth}|X}
    \textbf{Obiettivo} & \textbf{Descrizione} & \textbf{Metriche}  \\
    \hline
    Verifica &  &  \\
    \hline
    Gestione della qualità &  &  \\
    \end{tabularx}}

\end{comment}    
 % PROCESSI ORGANIZZATIVI   
\subsubsection{Processi organizzativi}
\begin{comment}
    {\renewcommand{\arraystretch}{1.5}
    \begin{tabularx}{\textwidth}{p{0.18\textwidth}|p{0.6\textwidth}|X}
    \textbf{Obiettivo} & \textbf{Descrizione} & \textbf{Metriche}  \\
    \hline
    Gestione organizzativa &  &  \\
    \end{tabularx}}
    \end{comment}
    


    

\begin{comment}
    {\renewcommand{\arraystretch}{1.5}
    \begin{tabularx}{\textwidth}{p{0.05\textwidth}|p{0.35\textwidth}|X|X}
    \textbf{ID} & \textbf{Nome metrica} & \textbf{Valore accettabile} & \textbf{Valore ottimale}  \\
    \hline
    \multicolumn{4}{l}{\cellcolor{primarycolor}\textbf{\textit{Gestione organizzativa}}} \\
    \hline
     &  &  &  \\
    \end{tabularx}}  
\end{comment}




%Test e specifiche
\newpage
\section{Verifica}


\subsection{Test di unità}
    {\renewcommand{\arraystretch}{1.5}
    \begin{tabularx}{\textwidth}{p{0.12\textwidth}|p{0.7\textwidth}|X}
    \textbf{ID} & \textbf{Descrizione} & \textbf{Stato}  \\
    \hline
     &  & \\
    \hline
     &  &  \\
    \hline
     &  & \\
    \hline
     &  &  \\
    \end{tabularx}}
    \subsubsection{Tracciamento}
    
\subsection{Test di integrazione}
    {\renewcommand{\arraystretch}{1.5}
    \begin{tabularx}{\textwidth}{p{0.12\textwidth}|p{0.7\textwidth}|X}
    \textbf{ID} & \textbf{Descrizione} & \textbf{Stato}  \\
    \hline
     &  & \\
    \hline
     &  &  \\
    \hline
     &  & \\
    \end{tabularx}}
    \subsubsection{Tracciamento}

\subsection{Test di sistema}
    {\renewcommand{\arraystretch}{1.5}
    \begin{tabularx}{\textwidth}{p{0.12\textwidth}|p{0.55\textwidth}|p{0.12\textwidth}|X}
    \textbf{ID} & \textbf{Descrizione} & \textbf{Requisito} & \textbf{Stato}  \\
    \hline
     &  & \\
    \hline
     &  &  \\
    \end{tabularx}}
    \subsubsection{Tracciamento}

\subsection{Test di accettazione}
    {\renewcommand{\arraystretch}{1.5}
    \begin{tabularx}{\textwidth}{p{0.12\textwidth}|p{0.7\textwidth}|X}
    \textbf{ID} & \textbf{Descrizione} & \textbf{Stato}  \\
    \hline
     &  & \\
    \hline
     &  &  \\
    \end{tabularx}}
    \subsubsection{Tracciamento}

\subsection{Checklist di controllo}

\subsubsection{Documentazione}
\subsubsection{Analisi dei Requisiti}
\subsubsection{Codifica}




\section{Dashboard metriche impiegate}
\subsection{Processi primari}
    
   
\subsection{Processi di supporto}
\begin{comment}
    {\renewcommand{\arraystretch}{1.5}
    \begin{tabularx}{\textwidth}{p{0.05\textwidth}|p{0.35\textwidth}|X|X}
    \textbf{ID} & \textbf{Nome metrica} & \textbf{Valore accettabile} & \textbf{Valore ottimale}  \\
    \hline
    \multicolumn{4}{l}{\cellcolor{primarycolor}\textbf{\textit{Verifica}}} \\
    \hline
     &  &  &  \\
    \hline
     &  &  &  \\
    \hline
     &  &  &  \\
    \hline
     &  &  &  \\
    \hline
    \multicolumn{4}{l}{\cellcolor{primarycolor}\textbf{\textit{Gestione della qualità}}} \\
    \hline
     &  &  &  \\
    \end{tabularx}}  
  
\subsubsection{Documentazione}
    {\renewcommand{\arraystretch}{1.5}
    \begin{tabularx}{\textwidth}{p{0.15\textwidth}|p{0.35\textwidth}|X|X}
    \textbf{ID} & \textbf{Nome metrica} & \textbf{Valore accettabile} & \textbf{Valore ottimale}  \\
    \hline
    \multicolumn{4}{l}{\cellcolor{primarycolor}\textbf{\textit{Funzionalità}}} \\
    \hline
    MPD-CR & Copertura dei requisiti & 100 & 100 \\
    
    \hline
   
    \multicolumn{4}{l}{\cellcolor{primarycolor}\textbf{\textit{Usabilità}}}
    \\
    \hline
     MPD-DIU & Difficoltà incontrate dall'utente & 2  & 0 \\
      
    \hline
   
    \multicolumn{4}{l}{\cellcolor{primarycolor}\textbf{\textit{Efficienza}}}
    \\
    \hline
     MPD-TA & Tempo di apprendimento & 10 minuti & 5 minuti \\
      \hline
     \multicolumn{4}{l}{\cellcolor{primarycolor}\textbf{\textit{Affidabiltà}}} \\
     MPD-TM & Tempo di risposta medio & 3 secondi & 2 secondi \\
      \hline
     \multicolumn{4}{l}{\cellcolor{primarycolor}\textbf{\textit{Portabiltà}}} \\
     MPD-VS & La capacità del programma di eseguire su più piattaforme & 100  & 100 \\
    \hline
     \multicolumn{4}{l}{\cellcolor{primarycolor}\textbf{\textit{Manutenibilità}}} \\
     MPD-CC & la capacità del software di essere modificato, 
     includendo correzioni, miglioramenti o adattamenti. & 30  & 45 \\
    \end{tabularx}}   
\end{comment} 
\subsection{Processi organizzativi}

\section{Valutazioni per il miglioramento}  
\begin{comment}
\section{Resoconto delle attività di verifica}

\subsection{Verifica della documentazione}
\subsubsection{Errori ortografici}
\subsubsection{Indici di Gulpease}

\subsection{Verifica dei processi}

\subsubsection{Estimated at completion}
\subsubsection{Budget variance e schedule variance}
\subsubsection{Actual cost e estimate to complete}
\subsubsection{Earned value e planned value}
\subsubsection{Requirements stability index e satisfied obligatory requirements}
\subsubsection{Code coverage back-end}
\subsubsection{Code coverage front-end}
\subsubsection{Passed test cases percentage}
\subsubsection{Failed test cases percentage}
\subsubsection{Comprensibilità del codice}
\end{comment}
\end{document}
